\chapter{Descrizione dello stage}
\label{chap:descrizione-stage}

\section{Introduzione al progetto}
Il modulo di previsione sviluppato durante lo stage stima la domanda futura di prodotti finiti per supportare la pianificazione \gls{mrp} di \textit{ErpAI}. Le stime si basano sullo storico delle fatture di vendita e vengono propagate lungo le distinte base per evidenziare i fabbisogni di semilavorati e materie prime. I risultati vengono messi a disposizione dei planner come proposte rivedibili, affiancando inizialmente il processo manuale.

Obiettivi principali:
\begin{itemize}
    \item ridurre \textit{stock-out} e fermate impianto anticipando materiali critici;
    \item limitare \textit{overstock} e costi di magazzino evitando produzioni o acquisti non necessari;
    \item fornire previsioni integrabili con il modulo di \textit{planning} e l’interfaccia utente per approvare o modificare le proposte.
\end{itemize}


\section{Analisi preventiva dei rischi}

Durante la fase iniziale sono stati individuati alcuni rischi specifici per il modulo di previsione e la sua integrazione:

\begin{risk}{Qualità e quantità dei dati storici}
    \riskdescription{dati di vendita incompleti o disomogenei possono ridurre l’accuratezza del modello}
    \risksolution{pulizia e normalizzazione dei dataset, regole di validazione condivise con i referenti di business}
    \label{risk:data-quality}
\end{risk}

\begin{risk}{Integrazione con l’\gls{erp} e le \gls{api}}
    \riskdescription{ritardi o incongruenze nello scambio dati tra modulo di previsione, \textit{planning} e \gls{erp}}
    \risksolution{contratto \gls{api} documentato (OpenAPI), ambiente di staging per test end-to-end prima del rilascio}
    \label{risk:integration-api}
\end{risk}

\begin{risk}{Performance e tempi di risposta}
    \riskdescription{tempi di addestramento o interrogazione troppo lunghi per l’uso operativo}
    \risksolution{limitazione delle feature non essenziali, caching dei risultati, pianificazione di job batch fuori orario}
    \label{risk:performance}
\end{risk}

\begin{risk}{Qualità di generalizzazione del modello}
    \riskdescription{rischio di \gls{overfitting} o \gls{underfitting} a causa di dati sporchi o parametrizzazione non ottimale}
    \risksolution{pulizia e normalizzazione dei dati, regolarizzazione e ottimizzazione degli iperparametri per bilanciare complessità e capacità predittiva}
    \label{risk:model-generalization}
\end{risk}

\begin{risk}{Adozione da parte dei planner}
    \riskdescription{diffidenza verso i risultati del modello e ritardi nell’uso in produzione}
    \risksolution{esecuzione parallela al processo manuale per più cicli, raccolta \textit{feedback} e correzioni iterative}
    \label{risk:user-adoption}
\end{risk}

\section{Pianificazione}
\subsection{Modalità di stage e interazione}
Lo stage si svolge in presenza, con allineamenti quotidiani e tre \textit{touchpoint} settimanali con il tutor aziendale per supporto tecnico e validazione degli step. La comunicazione avviene di persona; il tutor fornisce revisione del codice, risoluzione problemi e verifica delle scelte progettuali.

\subsection{Pianificazione settimanale}
La Tabella~\ref{tab:pianificazione_settimanale} riporta la pianificazione delle attività settimanali concordata con il tutor aziendale all’inizio dello stage.
\begin{table}[htbp]
    \centering
    \renewcommand{\arraystretch}{1.2}
    \rowcolors{1}{}{tableGray}
    \begin{tabular}{|c|p{8cm}|c|}
    \hline
    \textbf{Settimana} & \textbf{Obiettivi/Attività} & \textbf{Ore}\\ \hline
    1 & Analisi requisiti e rischi; studio delle tecnologie da utilizzare. & 40\\ \hline
    2 & Progettazione e revisione della documentazione esistente. & 40\\ \hline
    3 & Codifica del progetto. & 40\\ \hline
    4 & Codifica del progetto. & 40\\ \hline
    5 & Codifica del progetto. & 40\\ \hline
    6 & Codifica del progetto. & 40\\ \hline
    7 & Codifica dei test e stesura della documentazione. & 40\\ \hline
    8 & Preparazione della presentazione finale, test conclusivi e chiusura attività. & 20\\ \hline
    \end{tabular}
    \caption{Pianificazione settimanale delle attività di stage.}
    \label{tab:pianificazione_settimanale}
\end{table}
\newpage
